\chapter{Bandit Rate Adaptation with a Known Lower Bound on Service Slack}
\label{Section:ChapRateAdaptation}
\noindent\textbf{Source and author contributions.}
This chapter is based on the results for a known lower bound on service slack in
``Bandit-Based Rate Adaptation for a Single-Server Queue,'' joint work with
Mevan Wijewardena and Michael J. Neely~\cite{wijewardena2026banditbased}.
I developed the rate-adaptation method using this slack lower bound and was primarily
responsible for its proofs and presentation. Wijewardena developed the
original paper's lower-bound construction and proof. That result is
restated here as Theorem~\ref{ch5:lower:theorem} without proof. Neely contributed
primarily through advising and reviewing our work on the original paper.
The paper's adaptive algorithm for unknown slack is not included here.
The cited preprint is distributed under the
\href{https://creativecommons.org/licenses/by/4.0/}{Creative Commons Attribution
4.0 International license}. This chapter adapts its presentation and analysis.

The statement and proof of Theorem~\ref{ch5:upper:thm:ucb-log-lower} were
generated by AI for this dissertation and are not part of the original
joint paper.

\section{Introduction}
\label{ch5:sec:introduction}

Consider a transmitter that serves a queue over a channel whose capacity
is sampled independently each slot from an unknown, fixed distribution.
In each slot, it chooses a transmission rate without observing the current
channel capacity, then receives a binary acknowledgment indicating whether
the transmission succeeded.
A successful transmission supplies service to the queue. A failed
transmission supplies none. We seek a bound on time-average expected queue
length that holds uniformly over all horizons, including the learning period.

Chapter~\ref{Section:ChapControl}
studies feasible storage policies under full cost feedback. Here, the
transmitter must estimate the mean offered service of each grid rate from
transmission successes and failures, while new work continues to arrive.

The action space is the interval $[0,1]$. For a fixed rate, mean
\emph{offered service} is the rate multiplied by its probability of success.
It measures available service even when the queue contains insufficient work
to use it. Both the arrival distribution and the channel
distribution are unknown. We assume that the transmitter is given a positive
lower bound $\varepsilon$ on service slack, the difference between the best
achievable mean offered service and the mean arrival rate. This known slack
lower bound determines the rate grid.

\subsection{Related work}

Scheduling and resource allocation in networks have been studied with
objectives such as power efficiency, utility, and
fairness~\cite{Cruz2003,palomar2006tutorial,Songwu1999}. Lyapunov drift analysis is a
standard method for bounding time-average expected queue
length~\cite{neely2010stochastic}.

This chapter applies the classical UCB1 policy~\cite{Auer2002FinitetimeAO}
to a finite grid of transmission rates. After trying each rate once, the
policy chooses the rate with the largest estimated mean service plus an
exploration bonus. At a given time, this bonus is larger for rates tried
fewer times.
Background on bandit learning is given
in~\cite{Lai1985,cesa2006prediction,bubeck2012regret,lattimore2020bandit}.

In queueing systems, the controller must learn service performance while
jobs continue to arrive. Krishnasamy et al.~\cite{Krishnasamy2016,Krishnasamy2020}
study how learning unknown service probabilities increases the expected
queue length at each time, compared with a controller that knows those
probabilities. Fu and Modiano~\cite{Fu2023} study routing jobs among servers
to maximize the reward earned when jobs finish. Yang et al.~\cite{Yang2023}
study scheduling when service rates must be learned, including settings
where those rates change over time. Other work also considers changing
arrival and service conditions~\cite{Huang2024,Nguyen}.

We adapt the upper-bound analysis of Freund, Lykouris, and
Weng~\cite{Freund2023} to workloads in $[0,1]$ and service amounts determined
by the selected transmission rate. Their analysis bounds the increase in
time-average expected queue length caused by learning, relative to a
controller that knows the service characteristics, uniformly over all
horizons. Our bound controls the total time-average expected queue length.

Combes et al.~\cite{combes2019optimal} learn which transmission mode and
rate provide the highest throughput. Tang et al.~\cite{tang_joint_2021}
combine rate selection with detection of changes in channel conditions.
Other studies develop and evaluate bandit methods for wireless rate
control~\cite{cho_use_2023,tong_rate_2023,le_multi-armed_2024}.
Li et al.~\cite{li2020multi} survey related methods for allocating frequency
bands. Our model permits rates throughout a continuous interval, and our
performance criterion is time-average expected queue length.

Bandit learning with actions chosen from a continuous interval has been
studied under continuity assumptions on mean rewards~\cite{Agrawal1995}.
Other work constructs and refines finite sets of actions using bounds on
how much rewards can differ between nearby
actions~\cite{Kleinberg2008,pmlr-v134-podimata21a}.
Because lowering a rate preserves successful transmissions, we can use the
known slack lower bound to choose a grid containing a rate whose mean service
exceeds mean arrivals. This remains possible even when mean service is
discontinuous as a function of the attempted rate.

Learning to set prices provides a useful interpretation~\cite{Kleinberg2003}.
A purchase occurs when the buyer's value is at least the offered price,
just as a transmission succeeds when channel capacity is at least the
selected rate. Revenue corresponds to offered service. Rounding a price
down cannot prevent a purchase that would otherwise occur, and rounding a
rate down cannot prevent a successful transmission. Expected revenue or
mean offered service can decrease by at most the amount rounded off.
The pricing construction in~\cite{Kleinberg2003}
chooses the grid using the horizon. Here, the given lower bound on service
slack determines a grid that remains fixed for all horizons.

Audibert, Munos, and Szepesv{\'a}ri~\cite[Appendix~A.2,
Proposition~1 and Lemma~4]{audibert2009exploration} show that UCB1 continues
selecting a suboptimal arm even when both arms have deterministic
rewards. Their proof compares the two UCB indices. In our deterministic
example, the index comparison shows that UCB1 selects successful rates too
infrequently to keep up with arrivals over a chosen horizon. The resulting
accumulation of work gives the lower bound on time-average queue length.
\subsection{Contributions and scope}
\label{ch5:subsec:contributions}

We assume that arrivals are i.i.d., channel capacities are i.i.d., and the
two sequences are independent. Starting from an empty queue, UCB1 on a
suitably chosen finite rate grid achieves a time-average expected queue
bound of order
\[
O\!\left(\frac{1+\log(1/\varepsilon)}{\varepsilon^2}\right).
\]
This bound holds uniformly over all horizons.
Theorem~\ref{ch5:lower:theorem}, restated without proof
from~\cite{wijewardena2026banditbased}, shows
that every causal rate-selection policy faces some admissible channel and
horizon for which the time-average expected queue length is
$\Omega(\varepsilon^{-2})$. For the specific UCB1 policy studied here,
Theorem~\ref{ch5:upper:thm:ucb-log-lower} proves the stronger lower bound
$\Omega(\varepsilon^{-2}\log(1/\varepsilon))$, even with deterministic
arrivals and capacities. The logarithmic factor is therefore necessary
for this algorithm in the worst case. These bounds do not determine
whether another causal policy can remove this factor.



\section{Queueing model and performance criterion}
\label{ch5:sec:model}

We use the conventions in Section~\ref{sec:intro:notation}, with time
written as a parenthesized argument.

\subsection{Arrivals, service, and observations}

In a wireless link, modulation and coding determine how much data the
transmitter attempts to send in each slot. We model this choice using a
continuous range of rates. Workload is normalized so that the maximum
service per slot is 1. Arrivals are also bounded by 1 in the same units.

Time is divided into slots $t=1,2,\ldots$. Let $A(t)\in[0,1]$ denote the
arriving workload and $C(t)\in[0,1]$ the channel capacity in slot $t$.
The transmitter chooses an attempted rate $r(t)\in[0,1]$ without observing
$C(t)$. The resulting binary acknowledgment and offered service are
\begin{equation}
 B(t)\coloneqq\ind{r(t)\le C(t)},
 \qquad S(t)\coloneqq r(t)B(t).
 \label{ch5:eq:service}
\end{equation}
The transmitter observes $B(t)$ and computes $S(t)$ from its chosen rate.
Let $Q(t)$ be the queue length at the beginning of slot $t$, before that
slot's arrival and service. It satisfies
\begin{equation}
 Q(1)=0,\qquad
 Q(t+1)=[Q(t)+A(t)-S(t)]_+, \quad\forall t\ge1.
 \label{ch5:eq:queue}
\end{equation}
Arrivals can be served in the same slot. Actual departures equal
$\min\{Q(t)+A(t),S(t)\}$, whereas \eqref{ch5:eq:service} records offered
service. The acknowledgment is observed in every slot, even when the queue
is empty.

A causal policy can use the queue, its past attempted rates and
acknowledgments, and the arrivals observed so far. We allow the current
arrival $A(t)$ to be observed before the decision. Randomized decisions
use internal randomness independent of the arrivals and capacities.
The UCB policy developed below requires only observations through slot
$t-1$. It does not use $A(t)$ when choosing $r(t)$.

\begin{Thm:Assumption}[Stochastic environment]
\label{ch5:ass:environment}
The arrival sequence is i.i.d., the capacity sequence is i.i.d., and the
two sequences are independent. Their distributions are supported on
$[0,1]$ and are unknown to the transmitter. Write
$\lambda\coloneqq\E[A(1)]$.
\end{Thm:Assumption}

For a fixed transmission rate $r$, define its mean offered service by
\begin{equation}
 \mu(r)\coloneqq r\,\mathbb P(C(1)\ge r),\qquad r\in[0,1].
 \label{ch5:eq:g}
\end{equation}
For any causal selection of $r(t)$, independence of the current capacity
from the information used to select the rate gives
$\E[S(t)\mid r(t)]=\mu(r(t))$.
Equality $C(t)=r(t)$ counts as a successful transmission.
We use ``mean service'' as shorthand for mean offered service below.

\begin{Thm:Assumption}[Known lower bound on service slack]
\label{ch5:ass:slack}
There is a rate $r^*\in\argmax_{r\in[0,1]}\mu(r)$ and a given number
$\varepsilon>0$ such that
\begin{equation}
 \mu(r^*)\ge\lambda+\varepsilon.
 \label{ch5:eq:slack}
\end{equation}
The transmitter knows $\varepsilon$, but need not know $r^*$, $\mu(r^*)$,
or $\lambda$.
\end{Thm:Assumption}

Necessarily $0<\varepsilon\le1$: \eqref{ch5:eq:slack} gives
$\varepsilon\le \mu(r^*)-\lambda$, while \eqref{ch5:eq:g} gives $\mu(r^*)\le1$
and Assumption~\ref{ch5:ass:environment} gives $\lambda\ge0$.
The supplied $\varepsilon$ may be smaller than the actual slack
$\mu(r^*)-\lambda$. Supplying a smaller value of $\varepsilon$ can increase
the number of grid rates and increases the proved upper bound on
time-average expected queue length.

\subsection{Performance criterion and fixed-rate benchmark}

For an integer horizon $T\ge1$, let
\begin{equation}
 \overline Q_T\coloneqq\frac1T\sum_{t=1}^{T}\E[Q(t)].
 \label{ch5:eq:performance}
\end{equation}
Thus $\overline Q_T$ already includes expectation over the environment
and any controller randomization.
Our objective is to bound $\sup_{T\ge1}\overline Q_T$ using only
$\varepsilon$ and numerical constants, uniformly over the environments
satisfying Assumptions~\ref{ch5:ass:environment} and~\ref{ch5:ass:slack}.

The policy that always selects $r^*$ provides a fixed-rate benchmark
when the service characteristics are known.

\begin{Thm:Lemma}[Queue bound for the benchmark policy]
Under Assumptions~\ref{ch5:ass:environment} and~\ref{ch5:ass:slack},
the policy $r(t)=r^*$ for all $t$, started from $Q(1)=0$, satisfies
\begin{equation}
 \overline Q_T\le\frac{1}{2\varepsilon}
 \qquad\text{for every integer }T\ge1.
 \label{ch5:eq:oracle}
\end{equation}
\end{Thm:Lemma}

\begin{proof}
For $q\ge0$ and $d\in\mathbb R$,
\begin{equation}
 [q+d]_+^2-q^2\le(q+d)^2-q^2=2qd+d^2.
 \label{ch5:eq:positive-part-square}
\end{equation}
Applying \eqref{ch5:eq:positive-part-square}
to the queue update \eqref{ch5:eq:queue} gives
\begin{equation}
 Q(t+1)^2-Q(t)^2\le2Q(t)(A(t)-S(t))+(A(t)-S(t))^2.
 \label{ch5:eq:queue-square-increment}
\end{equation}
This inequality is samplewise and holds for every rate-selection policy.
For $r(t)=r^*$, Assumption~\ref{ch5:ass:environment} makes the current
arrival and service independent of $Q(t)$. Hence
\[
\begin{aligned}
 \E[A(t)-S(t)\mid Q(t)]
 &=\lambda-r^*\mathbb P(C(t)\ge r^*)\\
 &=\lambda-\mu(r^*)\le-\varepsilon,
 \qquad (A(t)-S(t))^2\le1.
\end{aligned}
\]
The equalities use \eqref{ch5:eq:service} and \eqref{ch5:eq:g}. The
inequalities use \eqref{ch5:eq:slack} and $A(t),S(t)\in[0,1]$.
Thus
\begin{align}
\E[Q(t+1)^2-Q(t)^2\mid Q(t)]
 &\overset{\text{(A)}}{\le}2Q(t)\E[A(t)-S(t)\mid Q(t)]+1\notag\\
 &\overset{\text{(B)}}{\le}1-2\varepsilon Q(t).
\label{ch5:eq:oracle-drift}
\end{align}
Here (A) uses \eqref{ch5:eq:queue-square-increment} and $(A(t)-S(t))^2\le1$.
Step (B) uses the conditional mean above.
Taking expectations, summing, and rearranging gives
\begin{align*}
2\varepsilon\sum_{t=1}^T\E[Q(t)]
 &\le T-\sum_{t=1}^T\E[Q(t+1)^2-Q(t)^2]\\
 &=T+\E[Q(1)^2]-\E[Q(T+1)^2]\le T.
\end{align*}
Here $Q(1)=0$. Division by $2\varepsilon T$ gives
\eqref{ch5:eq:oracle}.
\end{proof}

\section{Discretizing the rate space}
\label{ch5:sec:discretization}

Lowering a transmission rate by $d$ can reduce its mean offered service by
at most $d$. We use this fact to choose a finite grid containing at least
one rate whose mean service exceeds mean arrivals.

\begin{Thm:Lemma}[Mean-service bound for lowering a rate]
\label{ch5:lem:onesided}
For $0\le y\le x\le1$,
\begin{equation}
 \mu(y)\ge \mu(x)-(x-y).
 \label{ch5:eq:onesided}
\end{equation}
\end{Thm:Lemma}
\begin{proof}
By \eqref{ch5:eq:g},
\begin{align*}
\mu(y)
&=y\mathbb P(C(1)\ge y)
 \overset{\text{(A)}}{\ge}y\mathbb P(C(1)\ge x)\\
&=\mu(x)-(x-y)\mathbb P(C(1)\ge x)
 \ge\mu(x)-(x-y).
\end{align*}
Step (A) uses the inclusion
$\{C(1)\ge x\}\subseteq\{C(1)\ge y\}$.
\end{proof}

\begin{Thm:Remark}
There is no analogous bound for upward rounding. For example, if
$C(t)\equiv1/2$, then $\mu(1/2)=1/2$, whereas $\mu(r)=0$ for every
$r>1/2$. Thus, an arbitrarily small increase above rate $1/2$ causes
a mean-service loss of $1/2$.
\end{Thm:Remark}

For an integer $K\ge1$, use $K$ equally spaced positive rates:
\begin{equation}
 r_k\coloneqq\frac{k}{K},\qquad \mu_k\coloneqq\mu(r_k),\qquad k\in[K].
 \label{ch5:eq:grid-rates}
\end{equation}
Let
\begin{equation}
 k^*\in\argmax_{k\in[K]}\mu_k,
 \qquad \mu_K^*\coloneqq\mu_{k^*},\qquad \sigma\coloneqq\mu_K^*-\lambda.
 \label{ch5:eq:grid-slack}
\end{equation}
The actual grid slack $\sigma$ is used only in the analysis.

\begin{Thm:Lemma}[A grid retaining positive slack]
\label{ch5:lem:grid}
If $K\ge1/\varepsilon$, then
\[
 \mu_K^*\ge \mu(r^*)-\frac1K,
 \qquad \sigma\ge\varepsilon-\frac1K.
\]
In particular, choosing
\begin{equation}
 K\coloneqq\left\lceil\frac3\varepsilon\right\rceil
 \label{ch5:eq:d}
\end{equation}
ensures $\sigma\ge2\varepsilon/3>0$.
\end{Thm:Lemma}
\begin{proof}
By \eqref{ch5:eq:g}, \eqref{ch5:eq:slack}, $\lambda\ge0$, and
$K\ge1/\varepsilon$,
\[
 r^*\ge\mu(r^*)\ge\lambda+\varepsilon\ge\varepsilon\ge\frac1K.
\]
Thus
\begin{align}
 r^\downarrow&\coloneqq\frac{\lfloor Kr^*\rfloor}{K},
 \qquad \frac1K\le r^\downarrow\le r^*\le1,
 \label{ch5:eq:grid-rounded-rate}\\
 0&\le r^*-r^\downarrow<\frac1K.
 \label{ch5:eq:grid-rounding-distance}
\end{align}
Since $\lfloor Kr^*\rfloor\in\{1,\ldots,K\}$, $r^\downarrow$ is a
positive grid rate. Its mean service is at most $\mu_K^*$, so
\begin{align}
 \mu_K^*&\ge\mu(r^\downarrow)
 \overset{\mathrm{(A)}}{\ge}\mu(r^*)-(r^*-r^\downarrow)
 \ge\mu(r^*)-\frac1K,
 \label{ch5:eq:grid-mean-service-lower}\\
 \sigma&=\mu_K^*-\lambda\ge\mu(r^*)-\lambda-\frac1K
 \ge\varepsilon-\frac1K.
 \label{ch5:eq:grid-slack-lower}
\end{align}
Step (A) applies \eqref{ch5:eq:onesided}. The remaining bounds use
\eqref{ch5:eq:grid-rounding-distance} and \eqref{ch5:eq:slack}.
For the choice \eqref{ch5:eq:d},
\begin{align}
 K&=\left\lceil\frac3\varepsilon\right\rceil\ge\frac3\varepsilon
 \ge\frac1\varepsilon,\qquad \frac1K\le\frac\varepsilon3,
 \label{ch5:eq:grid-spacing-upper}\\
 \sigma&\ge\varepsilon-\frac1K\ge\frac{2\varepsilon}{3}>0.
 \label{ch5:eq:grid-positive-slack}
\end{align}
\end{proof}

\section{Rate selection on a fixed grid}
\label{ch5:upper:sec:algorithm}

Use the same grid of $K=\lceil3/\varepsilon\rceil$ rates from
Section~\ref{ch5:sec:discretization} for every horizon. Selecting a grid rate is now a
finite-armed bandit problem: each rate is an arm, and its offered service
is the observed reward. We apply UCB1~\cite{Auer2002FinitetimeAO} to these
arms.

Let $J(t)$ be the index of the rate selected in slot $t$.
For each rate, the controller records its selection count and average offered
service, with $S(t)=r_{J(t)}B(t)$. After slot $t$, these statistics are
\begin{equation}
 \begin{split}
 N_k(t)&\coloneqq\sum_{\tau=1}^{t}\ind{J(\tau)=k},\\
 \widehat\mu_k(t)&\coloneqq
 \frac{\sum_{\tau=1}^{t}S(\tau)\ind{J(\tau)=k}}
 {N_k(t)},\qquad N_k(t)>0.
 \end{split}
 \label{ch5:upper:eq:statistics}
\end{equation}
Define the confidence radius
\begin{equation}
 b(n,u)\coloneqq\sqrt{\frac{2\log u}{n}},\qquad
 b_k(t)\coloneqq b(N_k(t),t),
 \label{ch5:upper:eq:confidence-radius}
\end{equation}
where $n\ge1$, $u\ge1$, and the second expression requires $N_k(t)>0$.
The controller first samples every grid rate once. It then adds the
confidence radius to each estimated mean service to form the UCB index
\begin{equation}
 U_k(t)\coloneqq\widehat\mu_k(t)+b_k(t).
 \label{ch5:upper:eq:index}
\end{equation}
After initialization, it selects a rate in slot $t$ with the largest index
evaluated at $t-1$.

\begin{algorithm}[tbp]
\caption{UCB1 rate adaptation with a known lower bound on service slack}
\label{ch5:upper:alg:ucb}
\begin{algorithmic}[1]
\Require A known $\varepsilon>0$ satisfying Assumption~\ref{ch5:ass:slack}
\State Set $K\gets\lceil 3/\varepsilon\rceil$ and $r_k\gets k/K$ for $k=1,\ldots,K$
\State Initialize $N_k(0)\gets0$ for every $k$
\For{$t=1,2,\ldots$}
  \If{$t\le K$}
    \State $J(t)\gets t$
  \Else
    \State Choose $J(t)\in\argmax_{k\in[K]}U_k(t-1)$
    \State Resolve ties using a fixed rule
  \EndIf
  \State Set $r(t)\gets r_{J(t)}$, observe $B(t)$, and compute $S(t)\gets r(t)B(t)$
  \State Update the statistics in~\eqref{ch5:upper:eq:statistics}
\EndFor
\end{algorithmic}
\end{algorithm}

\section{Main guarantees}
\label{ch5:sec:guarantees}

The upper bounds control the queue under UCB1 uniformly over all horizons
and in the long run. One lower bound shows how large the time-average
expected queue length can be for any causal policy. The other gives a
larger worst-case lower bound specifically for UCB1.

\begin{Thm:Theorem}[Uniform bound on time-average expected queue length]
\label{ch5:upper:thm:main}
Under Assumptions~\ref{ch5:ass:environment} and~\ref{ch5:ass:slack},
let $0<\varepsilon\le1$ be the known slack lower bound and run
Algorithm~\ref{ch5:upper:alg:ucb} from $Q(1)=0$.
For every integer $T\ge1$,
\begin{equation}
 \overline Q_T
 \le \frac{170\bigl(1+\log(1/\varepsilon)\bigr)}{\varepsilon^2}.
 \label{ch5:upper:eq:uniform-bound}
\end{equation}
In particular, the worst-case time-average expected queue length is
$O\bigl(\varepsilon^{-2}\log(1/\varepsilon)\bigr)$ as $\varepsilon\downarrow0$.
\end{Thm:Theorem}

\begin{Thm:Remark}[Long-run queue length]
\label{ch5:upper:rem:long-run}
Under Assumptions~\ref{ch5:ass:environment} and~\ref{ch5:ass:slack}, run
Algorithm~\ref{ch5:upper:alg:ucb} from $Q(1)=0$.
For each fixed environment and supplied slack lower bound $\varepsilon$, the
long-run time-average expected queue length satisfies
\begin{equation}
 \limsup_{T\to\infty}\overline Q_T
 \le\frac3{4\varepsilon}.
 \label{ch5:upper:eq:long-run-bound}
\end{equation}
This long-run bound has the same $O(1/\varepsilon)$ order as the
fixed-rate benchmark bound~\eqref{ch5:eq:oracle}.
\end{Thm:Remark}

The following result is restated in this chapter's notation
from~\cite[Theorem~2]{wijewardena2026banditbased}.

\begin{Thm:Theorem}[A lower bound for every causal policy]
\label{ch5:lower:theorem}
For every $0<\varepsilon\le1/144$ and every causal, possibly randomized
rate-selection policy given $\varepsilon$, there is an environment
satisfying the model assumptions with $\mu(r^*)-\lambda=\varepsilon$ and
an integer horizon $T\ge1$ such that, starting from $Q(1)=0$,
\begin{equation}
 \overline Q_T
 \ge\frac{6\times10^{-7}}{\varepsilon^2}.
 \label{ch5:lower:main-bound}
\end{equation}
The arrivals in this environment are i.i.d.\ Bernoulli$(1/2)$ random
variables, independent of the i.i.d.\ channel capacities.
\end{Thm:Theorem}

The proof is given in~\cite[Appendix~J, p.~41 of the PDF]{wijewardena2026banditbased}
and is omitted here.

\begin{Thm:Theorem}[Necessity of the logarithmic factor for UCB1]
\label{ch5:upper:thm:ucb-log-lower}
There is a numerical constant $\varepsilon_0>0$ such that, for every
$0<\varepsilon\le\varepsilon_0$, there is an environment satisfying the
model assumptions with $\mu(r^*)-\lambda=\varepsilon$ in which
Algorithm~\ref{ch5:upper:alg:ucb}, started from $Q(1)=0$, satisfies
\begin{equation}
 \overline Q_{T_\varepsilon}
 \ge\frac{\log(1/\varepsilon)}{32\varepsilon^2}
 \label{ch5:upper:eq:ucb-log-lower}
\end{equation}
at a horizon $T_\varepsilon=\Theta(\varepsilon^{-3}\log(1/\varepsilon))$.
Thus, for UCB1, the upper bound in Theorem~\ref{ch5:upper:thm:main} and this lower bound have the same dependence on $\varepsilon$.
\end{Thm:Theorem}

The upper bound is proved in Section~\ref{ch5:upper:sec:bound}.
The UCB1 lower bound is proved in Section~\ref{ch5:lower:ucb-log-proof}
using deterministic arrivals and capacities.

\section{Upper-bound analysis}
\label{ch5:upper:sec:bound}

Recall that $K$ is the number of grid rates and
$\sigma=\mu_K^*-\lambda>0$ is the service slack of the best grid rate.
The grid slack satisfies
$\sigma\ge2\varepsilon/3$ by Lemma~\ref{ch5:lem:grid}.
For the analysis, define the mean-service gap of each grid rate and of the
selected rate by
\begin{equation}
 \Delta_k\coloneqq\mu_K^*-\mu_k,\quad k\in[K],
 \qquad \Delta(t)\coloneqq\Delta_{J(t)}.
 \label{ch5:upper:eq:slack-gaps}
\end{equation}
The controller does not need these gaps.

Let $\mathcal F_{t-1}$ contain the observations through slot $t-1$ and
the controller randomization used to select the rate in slot $t$.
This information determines both the UCB1 decision $J(t)$ and the
current queue length $Q(t)$. It excludes the current arrival $A(t)$
and capacity $C(t)$. Those two random variables are independent of
this history, so the selected gap $\Delta(t)$ gives the decomposition
\begin{equation}
 \E[A(t)-S(t)\mid\mathcal F_{t-1}]
 =\lambda-\mu_{J(t)}
 =-\sigma+\Delta(t).
 \label{ch5:upper:eq:conditional-mean}
\end{equation}
Following~\cite{Freund2023}, we use this identity to derive two bounds on
time-average expected queue length. Lemma~\ref{ch5:upper:lem:stage-one}
compares the actual queue with an auxiliary queue having a fixed negative
conditional mean net input. Lemma~\ref{ch5:upper:lem:stage-two} studies the
expected change in squared queue length and accounts for the mean service
lost by selecting suboptimal rates, weighted by the current queue length.
For fixed $K$ and $\sigma$, the first bound grows logarithmically with the
horizon, while the second approaches a constant as the horizon grows.
Section~\ref{ch5:upper:sec:combination}
therefore uses the first bound for shorter horizons and the second for
longer horizons to obtain a bound independent of $T$. This split is used
only in the proof. The controller follows the same UCB1 rule throughout.


\subsection{Accuracy of service estimates under adaptive sampling}
\label{ch5:upper:sec:concentration}

We use the standard concentration argument from the analysis of
UCB1~\cite{Auer2002FinitetimeAO}. For each rate, we apply Hoeffding's
inequality~\cite{hoeffding1963probability} at a fixed sample count,
then take a union over all possible counts to account for adaptive
sampling, as in~\cite[Lemma~3.5]{Freund2025transient}.

Here, $u$ denotes the number of completed slots and $n$ the number
of samples of a given rate. The decision in slot $t$ uses $u=t-1$.

We use the standard representation of a stochastic bandit by independent
sample sequences~\cite[Section~4.6]{lattimore2020bandit}. For each grid
rate $r_k$, generate independent service samples $X_{k,1},X_{k,2},\ldots$.
Each sample has the distribution of $r_k\ind{C(1)\ge r_k}$. On the $n$th
selection of that rate, reveal $X_{k,n}$. The sequences are independent
across rates and independent of the arrivals.

This construction does not change the distribution of what the controller
observes: given the past and the selected rate, the next service has
exactly the same distribution as in the original model. The observations,
decisions, and queue therefore have the same joint distribution as before.

Set $\widehat\mu_{k,n}\coloneqq n^{-1}\sum_{j=1}^{n}X_{k,j}$. In this construction,
\[
 \widehat\mu_k(t)=\widehat\mu_{k,N_k(t)},\qquad N_k(t)>0.
\]
Since $X_{k,j}\in[0,1]$, Hoeffding's inequality with the radius
in~\eqref{ch5:upper:eq:confidence-radius} gives, for each fixed sample
count $n\ge1$ and integer $u>1$,
\begin{equation}
\mathbb P\!\left(
  |\widehat\mu_{k,n}-\mu_k|\ge b(n,u)\right)
 \le2\exp\!\left(-2n b(n,u)^2\right)
=2\exp\!\left(-2n\frac{2\log u}{n}\right)
 =2u^{-4}.
\label{ch5:upper:eq:fixed-count}
\end{equation}

For each integer $u>1$, define
\begin{equation}
 \mathcal E_u
 \coloneqq\bigcap_{k=1}^{K}\bigcap_{n=1}^{u}
 \left\{\left|\widehat\mu_{k,n}-\mu_k\right|
 <b(n,u)\right\}.
 \label{ch5:upper:eq:confidence-event}
\end{equation}
A union bound in~\eqref{ch5:upper:eq:fixed-count} gives
\begin{equation}
 \mathbb P(\mathcal E_u^c)
 \le\sum_{k=1}^K\sum_{n=1}^u 2u^{-4}
 =\frac{2K}{u^3}.
 \label{ch5:upper:eq:confidence-envelope}
\end{equation}
On $\mathcal E_u$, each rate with $N_k(u)>0$ has
$1\le N_k(u)\le u$, so
\begin{equation}
 |\widehat\mu_k(u)-\mu_k|
 <b(N_k(u),u)=b_k(u).
 \label{ch5:upper:eq:adaptive-confidence}
\end{equation}
The inequality selects $n=N_k(u)$ in
\eqref{ch5:upper:eq:confidence-event}. Adding the bonus $b_k(u)$
to the resulting bounds on $\widehat\mu_k(u)$ and using
\eqref{ch5:upper:eq:index} yields
\begin{equation}
 \mu_k\le U_k(u)\le\mu_k+2b_k(u).
 \label{ch5:upper:eq:index-envelope}
\end{equation}

The next lemma is the classical UCB1 sampling bound of Auer,
Cesa-Bianchi, and Fischer~\cite[proof of Theorem~1]{Auer2002FinitetimeAO},
expressed in the notation of grid rates.

\begin{Thm:Lemma}[Sampling a suboptimal rate]
\label{ch5:upper:lem:counts}
For every grid rate with $\Delta_k>0$ and every integer $T\ge1$,
\begin{equation}
 \E[N_k(T)]\le\frac{8\log T}{\Delta_k^2}
 +1+\frac{\pi^2}{3}.
 \label{ch5:upper:eq:counts}
\end{equation}
\end{Thm:Lemma}

\begin{proof}
During initialization, each rate is sampled at most once, so
$N_k(T)\le1$ when $T\le K$. For $T>K$, set
\begin{equation}
m\coloneqq\left\lceil\frac{8\log T}{\Delta_k^2}\right\rceil.
\label{ch5:upper:eq:sample-threshold}
\end{equation}
For a decision at $K<t\le T$, put $u\coloneqq t-1$. If $N_k(u)\ge m$, then
\eqref{ch5:upper:eq:confidence-radius} gives
\begin{equation}
b_k(u)
=\sqrt{\frac{2\log u}{N_k(u)}}
 \overset{\text{(A)}}{\le}\sqrt{\frac{2\log T}{m}}
\overset{\text{(B)}}{\le}\sqrt{\frac{2\log T}{8\log T/\Delta_k^2}}
=\frac{\Delta_k}{2}.
\label{ch5:upper:eq:sample-radius-bound}
\end{equation}
Here, (A) uses $u<T$ and $N_k(u)\ge m$. Step (B) uses the ceiling lower
bound in \eqref{ch5:upper:eq:sample-threshold}.
For $K\le u<T$, define
\begin{equation}
\mathcal H_k(u)
 \coloneqq\{\widehat\mu_k(u)-\mu_k\ge b_k(u)\}
\cup\{\mu_K^*-\widehat\mu_{k^*}(u)\ge b_{k^*}(u)\}.
 \label{ch5:upper:eq:sample-error-event}
\end{equation}
By~\eqref{ch5:upper:eq:index} and~\eqref{ch5:upper:eq:slack-gaps},
if $N_k(u)\ge m$ and $\mathcal H_k(u)$ does not occur, the two indices satisfy
\begin{align}
U_k(u)
&=\widehat\mu_k(u)+b_k(u)
 \overset{\text{(A)}}{<}\mu_k+2b_k(u)
 \overset{\text{(B)}}{\le}\mu_k+\Delta_k=\mu_K^*,\notag\\
U_{k^*}(u)&=\widehat\mu_{k^*}(u)+b_{k^*}(u)
 \overset{\text{(C)}}{>}\mu_K^*.
\label{ch5:upper:eq:sample-index-comparison}
\end{align}
Steps (A) and (C) use the absence of the respective errors. Step (B) uses
\eqref{ch5:upper:eq:sample-radius-bound}.
Thus the optimal rate has a strictly larger index, and $r_k$ cannot be
selected. The probability of either error satisfies
\begin{align}
\mathbb P(\mathcal H_k(u))
 &\overset{\mathrm{(A)}}{\le}
 \sum_{n=1}^{u}
 \mathbb P(\widehat\mu_{k,n}-\mu_k\ge b(n,u))
+\sum_{n=1}^{u}
 \mathbb P(\mu_K^*-\widehat\mu_{k^*,n}\ge b(n,u))\notag\\
&\overset{\mathrm{(B)}}{\le}
 2\sum_{n=1}^{u}\exp(-2nb(n,u)^2)
 =2u^{-3}.
\label{ch5:upper:eq:sample-error-probability}
\end{align}
Here (A) takes a union over the possible sample counts $n=1,\ldots,u$.
Step (B) uses the one-sided Hoeffding bounds
underlying~\eqref{ch5:upper:eq:fixed-count} with the radius
\eqref{ch5:upper:eq:confidence-radius}.
On every sample path,
\begin{equation}
\begin{aligned}
 \{J(u+1)=k,\ N_k(u)\ge m\}&\subseteq\mathcal H_k(u),
 \qquad K\le u<T,\\
 N_k(T)&\le m+\sum_{u=K}^{T-1}\ind{\mathcal H_k(u)}.
\end{aligned}
\label{ch5:upper:eq:sample-path-count}
\end{equation}
The inclusion follows from~\eqref{ch5:upper:eq:sample-index-comparison}
and the maximizing rule in Algorithm~\ref{ch5:upper:alg:ucb}.
The count follows from~\eqref{ch5:upper:eq:statistics}: at most $m$
selections occur before the count reaches $m$, and the inclusion
applies to every later selection. Consequently,
\begin{equation*}
\E[N_k(T)]
\overset{\text{(A)}}{\le}m+2\sum_{u=K}^{T-1}u^{-3}
\overset{\text{(B)}}{\le}\frac{8\log T}{\Delta_k^2}+1
  +2\sum_{u=1}^{\infty}u^{-2}
=\frac{8\log T}{\Delta_k^2}+1+\frac{\pi^2}{3}.
\end{equation*}
Step (A) takes expectations in~\eqref{ch5:upper:eq:sample-path-count}
and applies~\eqref{ch5:upper:eq:sample-error-probability}. Step (B) uses
$\lceil z\rceil\le z+1$ in \eqref{ch5:upper:eq:sample-threshold},
$u^{-3}\le u^{-2}$ for $u\ge1$, and enlarges the summation range.
The last equality uses $\sum_{u=1}^{\infty}u^{-2}=\pi^2/6$.
\end{proof}

\subsection{Auxiliary-queue comparison bound}
\label{ch5:upper:sec:stage-one}

Fix $\delta\in(0,\sigma)$. Following~\cite{Freund2023}, we compare the
actual queue with an auxiliary queue whose conditional mean net input
is $-(\sigma-\delta)$.

\begin{Thm:Lemma}[Auxiliary-queue comparison bound]
\label{ch5:upper:lem:stage-one}
For every integer $T\ge2$ and every $\delta\in(0,\sigma)$,
\begin{equation}
 \overline Q_T
 \le\frac{1}{4(\sigma-\delta)}+\frac{\sigma-\delta}{2}
 +\frac{2K\log T}{\delta}.
 \label{ch5:upper:eq:stage-one}
\end{equation}
It follows that
\begin{equation}
 \overline Q_T
 \le\frac12+\frac1\sigma
 \left(\frac12+\sqrt{2K\log T}\right)^2.
 \label{ch5:upper:eq:stage-one-simple}
\end{equation}
\end{Thm:Lemma}

\begin{proof}
Define
\begin{equation}
 \zeta(t)\coloneqq(A(t)-\lambda)-(S(t)-\mu_{J(t)}).
 \label{ch5:upper:eq:centered-noise}
\end{equation}
Given the history $\mathcal F_{t-1}$, the selected rate is fixed. The
arrival and offered service are then independent, and each lies in $[0,1]$.
For any such random variable $X$,
\begin{equation}
\begin{aligned}
\operatorname{Var}(X\mid\mathcal F_{t-1})
 &=\E[(X-1/2)^2\mid\mathcal F_{t-1}]
-(\E[X\mid\mathcal F_{t-1}]-1/2)^2\\
&\le\E[(X-1/2)^2\mid\mathcal F_{t-1}]
 \le\frac14.
\end{aligned}
\label{ch5:upper:eq:bounded-conditional-variance}
\end{equation}
We have
\begin{equation}
 \begin{aligned}
\E[\zeta(t)\mid\mathcal F_{t-1}]&=0,\\
\E[\zeta(t)^2\mid\mathcal F_{t-1}]
 &=\operatorname{Var}(A(t)\mid\mathcal F_{t-1})
   +\operatorname{Var}(S(t)\mid\mathcal F_{t-1})
\le\frac14+\frac14=\frac12.
\end{aligned}
 \label{ch5:upper:eq:noise-moments}
\end{equation}
The mean identity follows from~\eqref{ch5:upper:eq:centered-noise}
and~\eqref{ch5:upper:eq:conditional-mean}. Conditional independence gives
the sum of variances, and~\eqref{ch5:upper:eq:bounded-conditional-variance}
bounds each variance by $1/4$.
The net input can be written as
\begin{equation}
A(t)-S(t)
 \overset{\mathrm{(A)}}{=}\zeta(t)+\lambda-\mu_{J(t)}
\overset{\mathrm{(B)}}{=}\zeta(t)-\sigma+\Delta(t).
\label{ch5:upper:eq:net-input-decomposition}
\end{equation}
Here (A) uses~\eqref{ch5:upper:eq:centered-noise}. Step (B)
uses~\eqref{ch5:eq:grid-slack} and~\eqref{ch5:upper:eq:slack-gaps}.
Set
\begin{equation}
 \begin{aligned}
 Q(t+1)&=[Q(t)+\zeta(t)-\sigma+\Delta(t)]_+,\\
 \widetilde Q(t+1)
 &\coloneqq[\widetilde Q(t)+\zeta(t)-(\sigma-\delta)]_+,
 \qquad \widetilde Q(1)\coloneqq0.
 \end{aligned}
 \label{ch5:upper:eq:coupled-queues}
\end{equation}

Put $\kappa\coloneqq\sigma-\delta>0$. Since $\mathcal F_{t-1}$
determines $\widetilde Q(t)$,
\begin{align}
&\E[\widetilde Q(t+1)^2-\widetilde Q(t)^2
       \mid\mathcal F_{t-1}]\notag\\
 &\quad\overset{\mathrm{(A)}}{\le}
 2\widetilde Q(t)\E[\zeta(t)-\kappa\mid\mathcal F_{t-1}]
+\E[(\zeta(t)-\kappa)^2\mid\mathcal F_{t-1}]\notag\\
 &\quad=2\widetilde Q(t)(\E[\zeta(t)\mid\mathcal F_{t-1}]-\kappa)
+\E[\zeta(t)^2\mid\mathcal F_{t-1}]\notag\\
 &\qquad-2\kappa\E[\zeta(t)\mid\mathcal F_{t-1}]+\kappa^2
 \label{ch5:upper:eq:aux-increment-moment}\\
 &\quad\overset{\mathrm{(B)}}{=}
 \E[\zeta(t)^2\mid\mathcal F_{t-1}]+\kappa^2-2\kappa\widetilde Q(t)\notag\\
 &\quad\overset{\mathrm{(C)}}{\le}
 \frac12+\kappa^2-2\kappa\widetilde Q(t).
\label{ch5:upper:eq:aux-one-slot-drift}
\end{align}
Step (A) applies~\eqref{ch5:eq:positive-part-square} to
\eqref{ch5:upper:eq:coupled-queues}. After the expansion in
\eqref{ch5:upper:eq:aux-increment-moment}, (B) and (C) use
\eqref{ch5:upper:eq:noise-moments}. Taking expectations and summing gives
\begin{align}
2\kappa\sum_{t=1}^{T}\E[\widetilde Q(t)]
 &\le T\left(\frac12+\kappa^2\right)
 -\sum_{t=1}^T\E[\widetilde Q(t+1)^2-\widetilde Q(t)^2]\notag\\
 &=T\left(\frac12+\kappa^2\right)
 +\E[\widetilde Q(1)^2]-\E[\widetilde Q(T+1)^2]\notag\\
 &\le T\left(\frac12+\kappa^2\right).
 \label{ch5:upper:eq:aux-summed-drift}
\end{align}
Here $\widetilde Q(1)=0$. Division by $2\kappa T$ gives
\begin{equation}
 \frac1T\sum_{t=1}^{T}\E[\widetilde Q(t)]
 \le\frac{1/2+\kappa^2}{2\kappa}
 =\frac{1}{4\kappa}+\frac\kappa2
 =\frac{1}{4(\sigma-\delta)}+\frac{\sigma-\delta}{2}.
 \label{ch5:upper:eq:aux-drift}
\end{equation}

Define
\begin{equation}
 L_\delta(t)\coloneqq\sum_{\tau=1}^{t}[\Delta(\tau)-\delta]_+
 =\sum_{k:\Delta_k>\delta}(\Delta_k-\delta)N_k(t),
 \qquad L_\delta(0)\coloneqq0.
 \label{ch5:upper:eq:excess-loss}
\end{equation}
We prove $Q(t)\le\widetilde Q(t)+L_\delta(t-1)$ by induction.
At $t=1$, all three quantities are zero. Assuming the inequality holds
at slot $t$, the coupled updates give
\begin{equation*}
\begin{aligned}
Q(t+1)
&\overset{\mathrm{(A)}}{\le}
 [\widetilde Q(t)+\zeta(t)-(\sigma-\delta)
   +L_\delta(t-1)+\Delta(t)-\delta]_+\\
&\le
 [\widetilde Q(t)+\zeta(t)-(\sigma-\delta)]_+
 +L_\delta(t-1)+[\Delta(t)-\delta]_+\\
&\overset{\mathrm{(B)}}{=}\widetilde Q(t+1)+L_\delta(t).
\end{aligned}
\end{equation*}
Step (A) uses the actual queue update in~\eqref{ch5:upper:eq:coupled-queues}
and the induction hypothesis.
Step (B) applies the auxiliary update in~\eqref{ch5:upper:eq:coupled-queues}
and the definition~\eqref{ch5:upper:eq:excess-loss}.
Thus the induction proves, for every $t\ge1$,
\begin{equation}
 Q(t)\le\widetilde Q(t)+\sum_{\tau=1}^{t-1}[\Delta(\tau)-\delta]_+
 =\widetilde Q(t)+L_\delta(t-1).
 \label{ch5:upper:eq:comparison}
\end{equation}
Taking expectations in~\eqref{ch5:upper:eq:comparison}, averaging over
$t=1,\ldots,T$, and applying~\eqref{ch5:upper:eq:aux-drift} gives
\begin{equation}
 \overline Q_T=\frac1T\sum_{t=1}^{T}\E[Q(t)]
 \le\frac{1}{4(\sigma-\delta)}+\frac{\sigma-\delta}{2}
 +\frac1T\sum_{\tau=0}^{T-1}\E[L_\delta(\tau)].
 \label{ch5:upper:eq:comparison-transfer}
\end{equation}
So far, we have used no property specific to UCB1. The comparison applies
to any rule that selects a grid rate using only the history before the
current arrival and capacity are drawn. We now use the UCB1 sampling
bound to control $L_\delta$.

For $\Delta>\delta>0$,
\begin{equation}
 \frac1{4\delta}-\frac{\Delta-\delta}{\Delta^2}
 =\frac{(\Delta-2\delta)^2}{4\delta\Delta^2}\ge0.
 \label{ch5:upper:eq:gap-product-bound}
\end{equation}

We also have
\begin{equation}
 \sum_{\tau=1}^{T-1}\log\tau
 \le\int_1^T\log x\,\mathrm dx=T\log T-T+1.
 \label{ch5:upper:eq:log-sum-bound}
\end{equation}
For $0<\Delta_k\le1$, the sampling bound~\eqref{ch5:upper:eq:counts}
and $1+\pi^2/3<8$ imply
\[
 \E[N_k(\tau)]
 \le\frac{8\log\tau}{\Delta_k^2}+1+\frac{\pi^2}{3}
 \le\frac{8(1+\log\tau)}{\Delta_k^2}.
\]
The upper bound $\Delta_k\le1$ follows from
$0\le\mu_k\le\mu_K^*\le1$. For each grid rate with $\Delta_k>\delta$,
\begin{equation}
\begin{aligned}
(\Delta_k-\delta)\sum_{\tau=1}^{T-1}\E[N_k(\tau)]
 &\overset{\mathrm{(A)}}{\le}
 8\frac{\Delta_k-\delta}{\Delta_k^2}
 \sum_{\tau=1}^{T-1}(1+\log\tau)\\
 &\overset{\mathrm{(B)}}{\le}\frac2\delta
 \sum_{\tau=1}^{T-1}(1+\log\tau)
 \overset{\mathrm{(C)}}{\le}\frac2\delta T\log T.
\end{aligned}
\label{ch5:upper:eq:single-rate-excess-sum}
\end{equation}
Here (A) sums the preceding sampling estimate over $\tau=1,\ldots,T-1$.
Step (B) uses~\eqref{ch5:upper:eq:gap-product-bound}. Step (C) uses
\eqref{ch5:upper:eq:log-sum-bound}.
Combining these estimates gives
\begin{align*}
\overline Q_T
&\overset{\mathrm{(A)}}{\le}\frac{1}{4(\sigma-\delta)}+\frac{\sigma-\delta}{2}
+\frac1T\sum_{k:\Delta_k>\delta}
(\Delta_k-\delta)\sum_{\tau=1}^{T-1}\E[N_k(\tau)]\\
&\overset{\mathrm{(B)}}{\le}\frac{1}{4(\sigma-\delta)}+\frac{\sigma-\delta}{2}
+\frac{2K\log T}{\delta}.
\end{align*}
Step (A) is~\eqref{ch5:upper:eq:comparison-transfer} after expanding
$L_\delta$ by~\eqref{ch5:upper:eq:excess-loss}. Step (B) applies
\eqref{ch5:upper:eq:single-rate-excess-sum} to at most $K$ rates.
This proves~\eqref{ch5:upper:eq:stage-one}. Since
$\sigma-\delta<\sigma=\mu_K^*-\lambda\le1$, its middle term is at most $1/2$.
To obtain~\eqref{ch5:upper:eq:stage-one-simple}, set
$R\coloneqq\sqrt{2K\log T}>0$ for this calculation and choose
\[
\delta\coloneqq\frac{2\sigma R}{1+2R}\in(0,\sigma).
\]
Then $\sigma-\delta=\sigma/(1+2R)$ and $2K\log T=R^2$, so
\begin{align}
\overline Q_T
 &\overset{\mathrm{(A)}}{\le}\frac12+\frac{1}{4(\sigma-\delta)}
   +\frac{R^2}{\delta}
=\frac12+\frac{1+2R}{4\sigma}
   +\frac{R(1+2R)}{2\sigma}\notag\\
&=\frac12+\frac{(1+2R)^2}{4\sigma}
=\frac12+\frac1\sigma
   \left(\frac12+\sqrt{2K\log T}\right)^2.
\label{ch5:upper:eq:combination-comparison}
\end{align}
Here (A) uses~\eqref{ch5:upper:eq:stage-one} and the preceding bound
$(\sigma-\delta)/2\le1/2$.
\end{proof}

\subsection{Queue-weighted drift bound}
\label{ch5:upper:sec:stage-two}

To bound $\E[Q(t)\Delta(t)]$, define
\begin{equation}
 h(t)\coloneqq
 \begin{cases}
  \Delta(t),&1\le t\le K,\\[1mm]
  2b_{J(t)}(t-1),&t>K.
 \end{cases}
 \label{ch5:upper:eq:radius-budget}
\end{equation}
We also use
\begin{equation}
 \Gamma(T)\coloneqq\frac13+8(1+\log T)\log T.
 \label{ch5:upper:eq:budget-function}
\end{equation}
Let
\[
 Z_T\coloneqq\sum_{t=1}^{T}Q(t).
\]
Then $\E[Z_T]=T\overline Q_T$ by~\eqref{ch5:eq:performance}.
Because the queue can grow by at most one unit per slot, a large peak
requires substantial backlog in the preceding slots. We use this
observation to bound the largest queue length in terms of the sum of queue
lengths. This argument follows~\cite[Lemma~3.4]{Freund2025transient} and
refines an inequality of Huang, Golubchik, and Huang~\cite{Huang2024}.

\begin{Thm:Lemma}[Two deterministic bounds]
\label{ch5:upper:lem:deterministic}
For every integer $T\ge1$, on every sample path,
\begin{equation}
 \sum_{t=1}^{T}h(t)^2\le K\Gamma(T),
 \qquad
 \sum_{t=1}^{T}Q(t)^2
 \le\sqrt2\left(\sum_{t=1}^{T}Q(t)\right)^{3/2}.
 \label{ch5:upper:eq:deterministic}
\end{equation}
\end{Thm:Lemma}

\begin{proof}
During initialization, each rate is
sampled once. Since $\Delta_k=\mu_K^*-\mu_k$ and $\mu_k=\mu(k/K)$ by
\eqref{ch5:upper:eq:slack-gaps} and~\eqref{ch5:eq:grid-rates}, respectively,
\begin{equation}
\begin{aligned}
\Delta_k&=\mu(k^*/K)-\mu(k/K)
  \overset{\mathrm{(A)}}{\le}\frac{k^*-k}{K}, &&k\le k^*,\\
\Delta_k&=\mu_K^*-\mu_k
  \overset{\mathrm{(B)}}{\le}\frac{k^*}{K}, &&k>k^*.
\end{aligned}
\label{ch5:upper:eq:initial-gap-bounds}
\end{equation}
Step (A) applies~\eqref{ch5:eq:onesided} with $x=k^*/K$ and
$y=k/K$. Step (B) uses $\mu_k\ge0$ and
$\mu_K^*=\mu(k^*/K)=(k^*/K)\mathbb P(C(1)\ge k^*/K)\le k^*/K$
from the service definition~\eqref{ch5:eq:g}.
Thus
\begin{equation}
\begin{aligned}
\sum_{k=1}^{K}\Delta_k^2
 &\overset{\mathrm{(A)}}{\le}\frac1{K^2}\sum_{j=0}^{k^*-1}j^2
 +(K-k^*)\frac{(k^*)^2}{K^2}\\
&=\frac{k^*(k^*-1)(2k^*-1)}{6K^2}
 +(K-k^*)\frac{(k^*)^2}{K^2}\\
&\overset{\mathrm{(B)}}{\le}\frac{(k^*)^3}{3K^2}
 +(K-k^*)\frac{(k^*)^2}{K^2}\\
&=\frac{(k^*)^2}{K}-\frac{2(k^*)^3}{3K^2}
=\frac K3-\frac{(K-k^*)^2(K+2k^*)}{3K^2}
 \le\frac K3.
\end{aligned}
\label{ch5:upper:eq:initial-gap-budget}
\end{equation}
Here (A) squares the nonnegative bounds
in~\eqref{ch5:upper:eq:initial-gap-bounds} and reindexes the first
sum by $j=k^*-k$. Step (B) uses
$(k^*-1)(2k^*-1)\le2(k^*)^2$, since $k^*\ge1$.
The first equality uses $\sum_{j=0}^{m-1}j^2=m(m-1)(2m-1)/6$.
Equation~\eqref{ch5:upper:eq:initial-gap-budget} also bounds the contribution if the horizon ends before
initialization is complete. If $T>K$, each further selection of rate $k$
increases its sample count by one. Its counts just before those selections
are therefore exactly $1,\ldots,N_k(T)-1$. Using
$h(t)^2=8\log(t-1)/N_{J(t)}(t-1)$, from
\eqref{ch5:upper:eq:radius-budget} and
\eqref{ch5:upper:eq:confidence-radius}, gives
\begin{equation}
\begin{aligned}
\sum_{t=K+1}^{T}h(t)^2
 &\overset{\mathrm{(A)}}{\le}8\log T\sum_{k=1}^{K}\sum_{n=1}^{N_k(T)-1}\frac1n\\
 &\overset{\mathrm{(B)}}{\le}8K\log T\sum_{n=1}^{T-1}\frac1n
 \overset{\mathrm{(C)}}{\le}8K(1+\log T)\log T.
\end{aligned}
\label{ch5:upper:eq:post-initial-radius-budget}
\end{equation}
In (A), use $\log(t-1)\le\log T$. Step (B) extends each nonnegative sum using
$N_k(T)\le T$. Step (C) uses the harmonic-sum bound
$\sum_{n=1}^{T-1}n^{-1}\le1+\int_1^{T-1}x^{-1}\,\mathrm dx
\le1+\log T$ for $T\ge2$. Combining
\eqref{ch5:upper:eq:initial-gap-budget} and
\eqref{ch5:upper:eq:post-initial-radius-budget} with
\eqref{ch5:upper:eq:budget-function} proves the first
bound in~\eqref{ch5:upper:eq:deterministic}.
Write $Q_T^{\max}\coloneqq\max_{1\le t\le T}Q(t)$, and choose a slot
$t_{\max}$ at which the maximum is attained.
The queue recursion~\eqref{ch5:eq:queue} gives
$Q(t+1)=[Q(t)+A(t)-S(t)]_+\le Q(t)+1$,
since $0\le A(t)\le1$ and $S(t)\ge0$. Iterating this inequality
from $Q(1)=0$ gives
\begin{equation}
 0\le Q(t)\le t-1.
 \label{ch5:eq:pathwise}
\end{equation}
Applying~\eqref{ch5:eq:pathwise} at $t_{\max}$ and iterating the one-slot
bound backward, respectively, gives
\begin{equation}
Q_T^{\max}\le t_{\max}-1,\qquad
Q(t_{\max}-j)\ge Q_T^{\max}-j
\quad(0\le j<\lceil Q_T^{\max}\rceil).
 \label{ch5:upper:eq:queue-before-peak}
\end{equation}
Hence
\begin{equation}
Z_T\overset{\mathrm{(A)}}{\ge}
 \sum_{j=0}^{\lceil Q_T^{\max}\rceil-1}(Q_T^{\max}-j)
\overset{\mathrm{(B)}}{\ge}
 \int_0^{Q_T^{\max}}(Q_T^{\max}-x)\,\mathrm dx
 =\frac{(Q_T^{\max})^2}{2}.
 \label{ch5:upper:eq:peak-area}
\end{equation}
In (A), use~\eqref{ch5:upper:eq:queue-before-peak} and the
nonnegativity of the remaining queue lengths. All displayed slots exist
because $Q_T^{\max}\le t_{\max}-1$. For (B), the decreasing nonnegative function
$Q_T^{\max}-x$ on $[0,Q_T^{\max}]$ is at most $Q_T^{\max}-j$ on each interval
$[j,\min\{j+1,Q_T^{\max}\}]$, whose length is at most one.
Consequently,
\[
\sum_{t=1}^{T}Q(t)^2
\le Q_T^{\max}Z_T
\overset{\mathrm{(A)}}{\le}\sqrt2\,Z_T^{3/2}.
\]
Step (A) uses $Q_T^{\max}\le\sqrt{2Z_T}$ from~\eqref{ch5:upper:eq:peak-area}.
\end{proof}

\begin{Thm:Lemma}[Queue-weighted drift bound]
\label{ch5:upper:lem:stage-two}
For every integer $T\ge1$,
\begin{equation}
 \sigma\overline Q_T\le\frac12+\frac4T
 +\left(\frac{2K^2\Gamma(T)^2}{T\sigma^3}\right)^{1/4}
  (\sigma\overline Q_T)^{3/4}.
 \label{ch5:upper:eq:implicit-drift}
\end{equation}
Consequently,
\begin{equation}
 \overline Q_T\le\frac2\sigma+\frac{16}{T\sigma}
 +\frac{2K^2\Gamma(T)^2}{T\sigma^4}.
 \label{ch5:upper:eq:stage-two}
\end{equation}
\end{Thm:Lemma}

\begin{proof}
Since $Q(t)$ is determined by the history $\mathcal F_{t-1}$ and $|A(t)-S(t)|\le1$,
taking conditional expectations in \eqref{ch5:eq:queue-square-increment}
and using \eqref{ch5:upper:eq:conditional-mean} gives
\begin{align}
 \frac12\E[Q(t+1)^2-Q(t)^2\mid\mathcal F_{t-1}]
 &\le\frac12+Q(t)(-\sigma+\Delta(t)),\notag\\
 \frac12\E[Q(t+1)^2-Q(t)^2]
 &\le\frac12-\sigma\E[Q(t)]+\E[Q(t)\Delta(t)].
 \label{ch5:upper:eq:drift}
\end{align}
For $t\le K$, the gap term equals $\E[Q(t)h(t)]$
by~\eqref{ch5:upper:eq:radius-budget}. For $t>K$, we
put $u\coloneqq t-1$ and recall the event
$\mathcal E_u$ from~\eqref{ch5:upper:eq:confidence-event}.
The following comparison is the standard UCB optimism
argument~\cite[Section~7.1]{lattimore2020bandit}, applied here to mean service.
On $\mathcal E_u$,
\begin{equation}
\begin{aligned}
 \mu_K^*
 &\overset{\mathrm{(A)}}{\le}U_{k^*}(u)
 \overset{\mathrm{(B)}}{\le}U_{J(t)}(u)
 \overset{\mathrm{(C)}}{\le}\mu_{J(t)}+2b_{J(t)}(u),\\
 \Delta(t)
 &=\mu_K^*-\mu_{J(t)}\le2b_{J(t)}(u)=h(t).
\end{aligned}
 \label{ch5:upper:eq:good-event-gap}
\end{equation}
Here (A) and (C) use~\eqref{ch5:upper:eq:index-envelope}.
Step (B) holds because UCB1 selects a maximizing index.
The second line uses the definitions
in~\eqref{ch5:upper:eq:slack-gaps}
and~\eqref{ch5:upper:eq:radius-budget}.
We therefore obtain
\begin{equation}
\E[Q(t)\Delta(t)]
\overset{\mathrm{(A)}}{\le}\E[Q(t)h(t)\ind{\mathcal E_u}]
  +u\mathbb P(\mathcal E_u^c)
\overset{\mathrm{(B)}}{\le}\E[Q(t)h(t)]+2Ku^{-2}.
\label{ch5:upper:eq:weighted-gap-bound}
\end{equation}
In (A), split the expectation over $\mathcal E_u$ and its complement,
then use~\eqref{ch5:upper:eq:good-event-gap} on the first event and
$Q(t)\Delta(t)\le u$ on the second, using $\Delta(t)\le1$ and
$Q(t)\le t-1=u$ from~\eqref{ch5:eq:pathwise}. Step (B) applies
$\mathbb P(\mathcal E_u^c)\le2Ku^{-3}$ from
\eqref{ch5:upper:eq:confidence-envelope}.
The total contribution of the failure events is bounded by
\begin{equation}
 2K\sum_{u=K}^{\infty}u^{-2}
 \overset{\mathrm{(A)}}{\le}2K\left(\frac1{K^2}+\int_K^{\infty}x^{-2}\,\mathrm dx\right)
 =2+\frac2K\le4.
 \label{ch5:upper:eq:failure-budget}
\end{equation}
Step (A) keeps the $u=K$ term and bounds the remaining sum by the
integral of the decreasing function $x^{-2}$.
Summing~\eqref{ch5:upper:eq:drift} over time and applying
\eqref{ch5:upper:eq:weighted-gap-bound}--\eqref{ch5:upper:eq:failure-budget} gives
\begin{align}
\sigma\E\left[\sum_{t=1}^{T}Q(t)\right]
 &\le\frac T2+4+\E\left[\sum_{t=1}^{T}Q(t)h(t)\right]
+\frac12\E[Q(1)^2]-\frac12\E[Q(T+1)^2]\notag\\
 &\overset{\mathrm{(A)}}{\le}\frac T2+4+\E\left[\sum_{t=1}^{T}Q(t)h(t)\right].
 \label{ch5:upper:eq:summed-drift}
\end{align}
Step (A) uses $Q(1)=0$. For $T\le K$, there are no terms from $\mathcal E_u^c$.

Set
\begin{equation}
 y_T\coloneqq\sigma\overline Q_T=\frac{\sigma\E[Z_T]}{T},
 \qquad
 a_T\coloneqq\frac{2K^2\Gamma(T)^2}{T\sigma^3}.
 \label{ch5:upper:eq:normalized-drift-quantities}
\end{equation}
On every sample path,
\begin{equation}
\begin{aligned}
\sum_{t=1}^{T}Q(t)h(t)
 &\overset{\mathrm{(A)}}{\le}\left(\sum_{t=1}^{T}Q(t)^2\right)^{1/2}
     \left(\sum_{t=1}^{T}h(t)^2\right)^{1/2}\\
 &\overset{\mathrm{(B)}}{\le}2^{1/4}\sqrt{K\Gamma(T)}\,Z_T^{3/4}.
\end{aligned}
\label{ch5:upper:eq:learning-term-pathwise}
\end{equation}
Here (A) is Cauchy--Schwarz~\eqref{eq:intro:holder-young};
(B) uses~\eqref{ch5:upper:eq:deterministic}. Thus
\begin{align*}
y_T
 &\overset{\mathrm{(A)}}{\le}\frac12+\frac4T
   +\frac1T\E\left[\sum_{t=1}^{T}Q(t)h(t)\right]\\
 &\overset{\mathrm{(B)}}{\le}\frac12+\frac4T
   +\frac{2^{1/4}\sqrt{K\Gamma(T)}}{T}\E[Z_T^{3/4}]\\
 &\overset{\mathrm{(C)}}{\le}\frac12+\frac4T
   +\frac{2^{1/4}\sqrt{K\Gamma(T)}}{T}(\E[Z_T])^{3/4}\\
&=\frac12+\frac4T
   +\frac{2^{1/4}\sqrt{K\Gamma(T)}}{T}
       \left(\frac{Ty_T}{\sigma}\right)^{3/4}\\
&=\frac12+\frac4T
   +\left(\frac{2K^2\Gamma(T)^2}{T\sigma^3}\right)^{1/4}
       y_T^{3/4}
 =\frac12+\frac4T+a_T^{1/4}y_T^{3/4}\\
 &\overset{\mathrm{(D)}}{\le}\frac12+\frac4T
    +\frac14a_T+\frac34y_T.
\end{align*}
Step (A) divides~\eqref{ch5:upper:eq:summed-drift} by $T$;
(B) uses~\eqref{ch5:upper:eq:learning-term-pathwise}.
The equalities use~\eqref{ch5:upper:eq:normalized-drift-quantities}
and give~\eqref{ch5:upper:eq:implicit-drift}.
For (C) and (D), concavity of $z\mapsto z^{3/4}$ gives the tangent bound
\[
 z^{3/4}\le\frac34r^{-1/4}z+\frac14r^{3/4},\qquad z\ge0,\quad r>0.
\]
For (C), set $r=\E[Z_T]$ and take expectations with $z=Z_T$;
if $\E[Z_T]=0$, then $Z_T=0$ almost surely and (C) is immediate.
For (D), set $r=a_T$, $z=y_T$, and multiply by $a_T^{1/4}$ to obtain
$a_T^{1/4}y_T^{3/4}\le a_T/4+3y_T/4$; the case $a_T=0$ is immediate.
Rearranging and dividing by $\sigma>0$ gives
\begin{align}
 \frac14y_T&\le\frac12+\frac4T+\frac14a_T,
 \qquad y_T\le2+\frac{16}{T}+a_T,\notag\\
 \overline Q_T&=\frac{y_T}{\sigma}
 \le\frac2\sigma+\frac{16}{T\sigma}
       +\frac{2K^2\Gamma(T)^2}{T\sigma^4}.
 \label{ch5:upper:eq:combination-drift}
\end{align}
This is~\eqref{ch5:upper:eq:stage-two}.
\end{proof}

\begin{proof}[Proof of Remark~\ref{ch5:upper:rem:long-run}]
With $K$ and $\sigma$ fixed, the definitions
in~\eqref{ch5:upper:eq:normalized-drift-quantities}
and $\Gamma(T)=O((\log T)^2)$ give $a_T\to0$.
The explicit estimate~\eqref{ch5:upper:eq:stage-two} shows that $y_T$
stays bounded, so the learning term in~\eqref{ch5:upper:eq:implicit-drift}
vanishes. Taking the limit superior in~\eqref{ch5:upper:eq:implicit-drift} and using
$y_T=\sigma\overline Q_T$ and $\sigma\ge2\varepsilon/3$
from~\eqref{ch5:eq:grid-positive-slack} proves the bound.
\end{proof}

\subsection{Combining the two bounds}
\label{ch5:upper:sec:combination}

\begin{proof}[Proof of Theorem~\ref{ch5:upper:thm:main}]
For $T=1$, $\overline Q_1=Q(1)=0$. For $T\ge2$, set
$\ell_\varepsilon\coloneqq\log(1/\varepsilon)\ge0$ and choose
\begin{equation}
 T_{\mathrm{cut}}\coloneqq e^{29/2}\varepsilon^{-6}.
 \label{ch5:upper:eq:cutoff-definition}
\end{equation}
Apply \eqref{ch5:upper:eq:combination-comparison} for
$2\le T\le T_{\mathrm{cut}}$ and
\eqref{ch5:upper:eq:combination-drift} for $T\ge T_{\mathrm{cut}}$.
The numerical estimates verified in
Appendix~\ref{ch5:upper:sec:constant-verification} give, for $T\ge2$,
\begin{equation*}
\varepsilon^2\overline Q_T
 \overset{\mathrm{(A)}}{\le}
 \max\{169.903+64\ell_\varepsilon,167.959\}
\le170+64\ell_\varepsilon
 \le170(1+\ell_\varepsilon).
\end{equation*}
Here (A) uses~\eqref{ch5:upper:eq:short-horizon-bound}
and~\eqref{ch5:upper:eq:long-horizon-bound} on their respective ranges,
together with~\eqref{ch5:upper:eq:short-numerical-constant}
and~\eqref{ch5:upper:eq:long-numerical-constant}. Together with the case $T=1$,
this proves~\eqref{ch5:upper:eq:uniform-bound}.
\end{proof}



\section{Lower-bound analysis}
\label{ch5:lower:section}

Theorem~\ref{ch5:lower:theorem} gives the lower bound for every causal
policy. Its proof is in the original paper. Here we prove the stronger
bound for the prescribed UCB1 policy by constructing a deterministic
environment in which only six grid rates provide positive service. UCB1
spends many slots trying the remaining rates, which always fail, allowing
incoming work to accumulate.

\subsection{A matching lower bound for UCB1}
\label{ch5:lower:ucb-log-proof}

The deterministic construction and index comparison have a two-arm
precedent in Audibert, Munos, and Szepesv{\'a}ri~\cite[Appendix~A.2,
Proposition~1 and Lemma~4]{audibert2009exploration}.
The proof below applies this approach to the prescribed rate grid
and queueing objective.

\begin{proof}[Proof of Theorem~\ref{ch5:upper:thm:ucb-log-lower}]
Take deterministic arrivals and capacities
\begin{equation}
 A(t)\equiv\varepsilon,\qquad C(t)\equiv2\varepsilon,
 \qquad 0<\varepsilon\le\frac1{12}.
 \label{ch5:lower:ucb-environment}
\end{equation}
Here $\lambda=\varepsilon$ and $\mu(r)=r\ind{r\le2\varepsilon}$,
so $\mu(r^*)=2\varepsilon$ and $\mu(r^*)-\lambda=\varepsilon$.
For the prescribed grid $r_k=k/K$, the choice
$K=\lceil3/\varepsilon\rceil$ gives $6\le2\varepsilon K<7$.
Thus exactly the first six grid rates provide positive service.
Define
\begin{equation}
\begin{aligned}
 \mathcal K_+&\coloneqq\{1,\ldots,6\},\qquad
 \mathcal K_0\coloneqq\{7,\ldots,K\},\\
 K_0&\coloneqq|\mathcal K_0|=K-6,\\
 \mu_k&=r_k\ind{k\in\mathcal K_+},\qquad
 0\le\mu_k\le2\varepsilon\ind{k\in\mathcal K_+}.
\end{aligned}
\label{ch5:lower:ucb-arm-properties}
\end{equation}
Set
\begin{equation}
 T_\varepsilon\coloneqq
 \left\lfloor\frac{K_0\log K_0}{8\varepsilon^2}\right\rfloor,
 \qquad T\coloneqq T_\varepsilon.
 \label{ch5:lower:ucb-horizon}
\end{equation}
The grid choice~\eqref{ch5:eq:d} and~\eqref{ch5:lower:ucb-arm-properties}
give $K_0=3/\varepsilon+O(1)$. Hence~\eqref{ch5:lower:ucb-horizon} gives
\begin{equation}
 T=\Theta\!\left(\varepsilon^{-3}\log(1/\varepsilon)\right)
 \qquad(\varepsilon\downarrow0).
 \label{ch5:lower:ucb-horizon-scale}
\end{equation}
Choose $\varepsilon_0\in(0,1/12]$ so that, for every
$0<\varepsilon\le\varepsilon_0$,
\begin{equation}
\begin{gathered}
 \frac2\varepsilon\le K_0\le\frac3\varepsilon,\qquad K<T,\qquad
 \frac{K_0\log K_0}{8\varepsilon^2}\ge2,\\
 \frac{48}{K_0}+\frac{25}{2T}\le\frac14.
\end{gathered}
\label{ch5:lower:ucb-small-slack}
\end{equation}
The bounds on $K_0$ follow from $K_0=\lceil3/\varepsilon\rceil-6$
for $\varepsilon\le1/12$. The remaining conditions hold by
\eqref{ch5:lower:ucb-horizon-scale} and $K_0=3/\varepsilon+O(1)$.
Assume $0<\varepsilon\le\varepsilon_0$ for the remainder of the proof.

We first bound how often rates in $\mathcal K_+$ are selected by comparing
their indices with the index of a rate in $\mathcal K_0$.

Each observation of arm $k$ equals $\mu_k$. With the fixed tie-breaking
rule in Algorithm~\ref{ch5:upper:alg:ucb}, all selections, counts, and queue
lengths are therefore deterministic. After initialization, the statistics and
indices in~\eqref{ch5:upper:eq:statistics} and~\eqref{ch5:upper:eq:index}
specialize to
\begin{equation}
\begin{aligned}
 \widehat\mu_k(u)&=\mu_k,\qquad U_k(u)=\mu_k+b_k(u),\\
 J(u+1)&\in\argmax_{k\in[K]}U_k(u),\qquad u\ge K.
\end{aligned}
\label{ch5:lower:ucb-index-definitions}
\end{equation}
Fix $K\le u<T$ and choose $j\in\argmin_{\ell\in\mathcal K_0}N_\ell(u)$.
Then
\begin{align}
 N_j(u)&\le\frac1{K_0}\sum_{\ell\in\mathcal K_0}N_\ell(u)\le\frac{u}{K_0},
 \label{ch5:lower:ucb-light-count}\\
b_j(u)^2
 &\overset{\mathrm{(A)}}{\ge}\frac{2K_0\log u}{u}
\overset{\mathrm{(B)}}{\ge}\frac{2K_0\log K_0}{T}
 \overset{\mathrm{(C)}}{\ge}16\varepsilon^2.
\label{ch5:lower:ucb-light-bonus}
\end{align}
Step (A) uses~\eqref{ch5:lower:ucb-light-count} and~\eqref{ch5:upper:eq:confidence-radius};
(B) uses $K_0<K\le u<T$; (C) uses~\eqref{ch5:lower:ucb-horizon}.
If $J(u+1)=k\in\mathcal K_+$, then
\begin{equation}
b_j(u)=U_j(u)\overset{\mathrm{(A)}}{\le} U_k(u)=\mu_k+b_k(u)
\overset{\mathrm{(B)}}{\le}2\varepsilon+b_k(u)
\overset{\mathrm{(C)}}{\le}\tfrac12b_j(u)+b_k(u).
\label{ch5:lower:ucb-index-comparison}
\end{equation}
Step (A) uses~\eqref{ch5:lower:ucb-index-definitions};
(B) uses $\mu_k\le2\varepsilon$ from~\eqref{ch5:lower:ucb-arm-properties};
(C) uses~\eqref{ch5:lower:ucb-light-bonus}. By~\eqref{ch5:upper:eq:confidence-radius},
\begin{equation}
\begin{gathered}
 b_k(u)\ge\tfrac12b_j(u),\qquad
 \frac{2\log u}{N_k(u)}=b_k(u)^2
 \ge\frac{b_j(u)^2}{4}=\frac{2\log u}{4N_j(u)},\\
 N_k(u)\le4N_j(u)
 \overset{\mathrm{(A)}}{\le}\frac{4u}{K_0}.
\end{gathered}
 \label{ch5:lower:ucb-relative-count}
\end{equation}
Here $\log u>0$ and $N_k(u),N_j(u)>0$ since $u\ge K\ge3$;
(A) uses~\eqref{ch5:lower:ucb-light-count}.
For each $k\in\mathcal K_+$, apply~\eqref{ch5:lower:ucb-relative-count}
just before its last selection in slots $K+1,\ldots,T$ to obtain
\begin{equation}
 N_k(T)\le\frac{4T}{K_0}+1,\qquad k\in\mathcal K_+.
 \label{ch5:lower:ucb-per-arm-count}
\end{equation}
If there is no such selection, $N_k(T)=1$.
Define
\[
 N_+(t)\coloneqq\sum_{k\in\mathcal K_+}N_k(t),\qquad N_+(0)=0.
\]
By~\eqref{ch5:lower:ucb-per-arm-count},
\begin{equation}
 N_+(T)\le\frac{24T}{K_0}+6.
 \label{ch5:lower:ucb-event-count}
\end{equation}

The service and queue equations~\eqref{ch5:eq:service}
and~\eqref{ch5:eq:queue} specialize to
\begin{equation}
\begin{aligned}
 S(t)&=r_{J(t)}\ind{J(t)\in\mathcal K_+}
       \le2\varepsilon\ind{J(t)\in\mathcal K_+},\\
 Q(t+1)&=[Q(t)+\varepsilon-S(t)]_+.
\end{aligned}
\label{ch5:lower:ucb-queue-recursion}
\end{equation}
Using the service bound~\eqref{ch5:lower:ucb-arm-properties} and iterating
\eqref{ch5:lower:ucb-queue-recursion} from $Q(1)=0$ gives
\begin{equation}
 Q(t)\ge\varepsilon(t-1)-\sum_{\tau=1}^{t-1}S(\tau)
       \ge\varepsilon(t-1)-2\varepsilon N_+(t-1).
 \label{ch5:lower:ucb-queue-pointwise}
\end{equation}
Since the queue and counts are deterministic here,
$\overline Q_T=T^{-1}\sum_{\tau=1}^T Q(\tau)$. Therefore
\begin{align}
\overline Q_T
 &\overset{\mathrm{(A)}}{\ge}\frac{\varepsilon(T-1)}2
 -\frac{2\varepsilon}{T}\sum_{t=1}^{T}N_+(t-1)\notag\\
 &\overset{\mathrm{(B)}}{\ge}\frac{\varepsilon(T-1)}2-2\varepsilon N_+(T)\notag\\
 &\overset{\mathrm{(C)}}{\ge}\frac{\varepsilon(T-1)}2
       -2\varepsilon\left(\frac{24T}{K_0}+6\right)\notag\\
&=\varepsilon T\left(\frac12-\frac{48}{K_0}-\frac{25}{2T}\right)
 \overset{\mathrm{(D)}}{\ge}\frac{\varepsilon T}{4}.
\label{ch5:lower:ucb-queue-K-bound}
\end{align}
Step (A) sums~\eqref{ch5:lower:ucb-queue-pointwise};
(B) uses $N_+(t-1)\le N_+(T)$; (C) uses~\eqref{ch5:lower:ucb-event-count};
(D) uses~\eqref{ch5:lower:ucb-small-slack}.
Finally,
\begin{equation}
\overline Q_T
 \ge\frac{\varepsilon}{4}
       \frac{K_0\log K_0}{16\varepsilon^2}
 =\frac{K_0\log K_0}{64\varepsilon}
\ge\frac{\log(1/\varepsilon)}{32\varepsilon^2}.
\label{ch5:lower:ucb-final-bound}
\end{equation}
The first inequality uses \eqref{ch5:lower:ucb-queue-K-bound} and
$\lfloor z\rfloor\ge z/2$ for $z\ge2$, with
$z=K_0\log K_0/(8\varepsilon^2)\ge2$ by
\eqref{ch5:lower:ucb-small-slack}. The second uses
$K_0\ge2/\varepsilon$.
This proves the stated lower bound for every
$0<\varepsilon\le\varepsilon_0$.
\end{proof}

\section{Discussion and conclusion}
\label{ch5:sec:conclusion}

The controller needs only a known lower bound $\varepsilon$ on service slack to choose
a finite rate grid and run UCB1. Starting from an empty queue, it
achieves the time-average expected queue bound
$O((1+\log(1/\varepsilon))/\varepsilon^2)$ in
Theorem~\ref{ch5:upper:thm:main}. This bound holds for every horizon.

Theorem~\ref{ch5:lower:theorem}, proved in~\cite{wijewardena2026banditbased},
shows that, for every sufficiently small $\varepsilon$, any causal policy that knows
$\varepsilon$ but not the channel distribution faces a feasible
channel for which the time-average expected queue is
$\Omega(\varepsilon^{-2})$ at some horizon.
For the prescribed UCB1 policy,
Theorem~\ref{ch5:upper:thm:ucb-log-lower} gives the stronger bound
$\Omega(\varepsilon^{-2}\log(1/\varepsilon))$, even with deterministic
arrivals and capacities.

For each fixed environment and supplied $\varepsilon$, the long-run
bound~\eqref{ch5:upper:eq:long-run-bound} is $O(\varepsilon^{-1})$.
The bound that holds for every horizon must also cover queue buildup while
UCB1 learns which rates provide enough service. The UCB1 lower bound
exhibits such buildup at $T_\varepsilon$. The long-run bound concerns
time averages over increasingly long horizons; the contribution of any
fixed initial set of slots vanishes in these averages. It remains unresolved
whether another causal policy, under the same information assumptions, can
guarantee a time-average expected queue length of $O(\varepsilon^{-2})$ for
every horizon and every environment satisfying the model assumptions.

\clearpage
\begingroup
\renewcommand{\thesection}{\thechapter.\Alph{section}}
\renewcommand{\theHsection}{\theHchapter.appendix.\Alph{section}}
\setcounter{section}{0}
\section{Appendix: Verification of the numerical constant}
\label{ch5:upper:sec:constant-verification}

We verify the two numerical estimates used in the proof of
Theorem~\ref{ch5:upper:thm:main}. Set
$\ell_\varepsilon\coloneqq\log(1/\varepsilon)\ge0$ and
$c\coloneqq29/2$. The cutoff in~\eqref{ch5:upper:eq:cutoff-definition}
is therefore $T_{\mathrm{cut}}=e^c\varepsilon^{-6}$.

To account for rounding the grid size up to an integer, put
$\kappa\coloneqq K\varepsilon$. Since $K=\lceil3/\varepsilon\rceil$,
$3\le\kappa\le4$. Consequently,
\begin{align}
 \frac1\sigma
 &\overset{\text{(A)}}{\le}\frac1\varepsilon\frac{\kappa}{\kappa-1}
 \overset{\text{(B)}}{\le}\frac3{2\varepsilon},
 \label{ch5:upper:eq:slack-coefficient}\\
 \frac K\sigma
 &\overset{\text{(A)}}{\le}\frac1{\varepsilon^2}\frac{\kappa^2}{\kappa-1}
 \overset{\text{(B)}}{\le}\frac{16}{3\varepsilon^2},
 \label{ch5:upper:eq:linear-grid-coefficient}\\
 \frac{\sqrt{2K}}\sigma
 &\overset{\text{(A)}}{\le}\frac1{\varepsilon^{3/2}}\frac{\sqrt{2\kappa^3}}{\kappa-1}
 \overset{\text{(B)}}{\le}\sqrt{\frac{128}{9}}\,\varepsilon^{-3/2},
 \label{ch5:upper:eq:square-root-grid-coefficient}\\
 \frac{K^2}{\sigma^4}
 &\overset{\text{(A)}}{\le}\frac1{\varepsilon^6}\frac{\kappa^6}{(\kappa-1)^4}
 \overset{\text{(B)}}{\le}\frac{4096}{81\varepsilon^6}.
 \label{ch5:upper:eq:drift-grid-coefficient}
\end{align}
In these four bounds, (A) substitutes $K=\kappa/\varepsilon$ and
$\sigma\ge\varepsilon(\kappa-1)/\kappa$, the latter being
\eqref{ch5:eq:grid-slack-lower}. For (B), $\kappa/(\kappa-1)=1+1/(\kappa-1)$ is at most
$3/2$ because $\kappa\ge3$. The last three coefficients are largest at $\kappa=4$.
To check this, note
that on $[3,4]$,
\[
 \frac{\mathrm d}{\mathrm d\kappa}\frac{\kappa^2}{\kappa-1}
 =\frac{\kappa(\kappa-2)}{(\kappa-1)^2}\ge0,
 \qquad
 \frac{\mathrm d}{\mathrm d\kappa}\frac{\kappa^3}{(\kappa-1)^2}
 =\frac{\kappa^2(\kappa-3)}{(\kappa-1)^3}\ge0.
\]
The third coefficient is the square root of $2\kappa^3/(\kappa-1)^2$.
The fourth is $\kappa^6/(\kappa-1)^4=(\kappa^3/(\kappa-1)^2)^2$.

First consider the shorter horizons, $2\le T\le T_{\mathrm{cut}}$. The cutoff gives
\begin{equation}
 \begin{aligned}
 \log T&\overset{\text{(A)}}{\le}c+6\ell_\varepsilon,\\
 \varepsilon\log T
 &\le e^{-\ell_\varepsilon}
   (c+6\ell_\varepsilon)
 \overset{\text{(B)}}{\le}c.
 \end{aligned}
 \label{ch5:upper:eq:short-log-bounds}
\end{equation}
Here (A) takes logarithms of $T\le T_{\mathrm{cut}}=e^c\varepsilon^{-6}$
from~\eqref{ch5:upper:eq:cutoff-definition}.
Step (B) evaluates the decreasing
function at $\ell_\varepsilon=0$, since $c=29/2>6$ gives, for
$\ell_\varepsilon\ge0$,
\[
 \frac{\mathrm d}{\mathrm d\ell_\varepsilon}
 \left[e^{-\ell_\varepsilon}(c+6\ell_\varepsilon)\right]
 =e^{-\ell_\varepsilon}(6-c-6\ell_\varepsilon)<0.
\]
The constant terms in the comparison bound satisfy
\begin{equation}
 \frac{\varepsilon^2}{2}+\frac{\varepsilon^2}{4\sigma}
 \overset{\text{(A)}}{\le}\frac{\varepsilon^2}{2}+\frac{3\varepsilon}{8}
 \overset{\text{(B)}}{\le}\frac12+\frac38=\frac78,
 \label{ch5:upper:eq:short-constant-bound}
\end{equation}
where (A) uses \eqref{ch5:upper:eq:slack-coefficient} and (B) uses
$0<\varepsilon\le1$. We obtain
\begin{align}
\varepsilon^2\overline Q_T
 &\overset{\text{(A)}}{\le}\frac{\varepsilon^2}{2}+\frac{\varepsilon^2}{4\sigma}
 +\frac{\varepsilon^2\sqrt{2K}}{\sigma}\sqrt{\log T}
+\frac{2\varepsilon^2K}{\sigma}\log T\notag\\
&\overset{\text{(B)}}{\le}\frac78+\sqrt{\frac{128}{9}\,\varepsilon\log T}
 +\frac{32}{3}\log T\notag\\
&\overset{\text{(C)}}{\le}\frac78+\sqrt{\frac{128c}{9}}
 +\frac{32c}{3}+64\ell_\varepsilon.
\label{ch5:upper:eq:short-horizon-bound}
\end{align}
Step (A) uses \eqref{ch5:upper:eq:combination-comparison}. Step (B) uses
\eqref{ch5:upper:eq:short-constant-bound},
\eqref{ch5:upper:eq:square-root-grid-coefficient}, and
\eqref{ch5:upper:eq:linear-grid-coefficient} for the three terms,
respectively. Step (C) applies both bounds in
\eqref{ch5:upper:eq:short-log-bounds}.

Now consider the longer horizons, $T\ge T_{\mathrm{cut}}$, and
use~\eqref{ch5:upper:eq:combination-drift}. We will show that
its learning term decreases throughout this range. With $z\coloneqq\log T$, set
\[
 G(z)\coloneqq\Gamma(e^z)=\frac13+8z+8z^2.
\]
The factor $\Gamma(T)^2/T$ then becomes $e^{-z}G(z)^2$.
For $z>0$,
\begin{align}
 2G(z)-zG'(z)&=\frac23+8z>0,
 \qquad \frac{G'(z)}{G(z)}\le\frac2z,
 \label{ch5:upper:eq:G-log-derivative}\\
 \frac{\mathrm d}{\mathrm dz}\log\!\left(e^{-z}G(z)^2\right)
 &=-1+2\frac{G'(z)}{G(z)}\le-1+\frac4z.
 \label{ch5:upper:eq:horizon-monotonicity}
\end{align}
Thus $\Gamma(T)^2/T$ is nonincreasing for $T\ge e^4$.
Since $T_{\mathrm{cut}}\ge e^c>e^4$, substituting
$T_{\mathrm{cut}}=e^c\varepsilon^{-6}=e^{c+6\ell_\varepsilon}$ and
$\Gamma(T_{\mathrm{cut}})=G(c+6\ell_\varepsilon)$ gives
\begin{align}
 \overline Q_T
 &\overset{\text{(A)}}{\le}\frac3\varepsilon+\frac{24}{T\varepsilon}
 +\frac{8192}{81\varepsilon^6}\frac{\Gamma(T)^2}{T}\notag\\
 &\overset{\text{(B)}}{\le}\frac3\varepsilon+\frac{24}{T_{\mathrm{cut}}\varepsilon}
 +\frac{8192}{81\varepsilon^6}\frac{\Gamma(T_{\mathrm{cut}})^2}{T_{\mathrm{cut}}}\notag\\
 &=\frac3\varepsilon+24e^{-c}\varepsilon^5
 +\frac{8192}{81}e^{-c}G(c+6\ell_\varepsilon)^2.
 \label{ch5:upper:eq:long-cutoff-bound}
\end{align}
Here (A) applies \eqref{ch5:upper:eq:slack-coefficient} and
\eqref{ch5:upper:eq:drift-grid-coefficient} to
\eqref{ch5:upper:eq:combination-drift}. Step (B) uses $T\ge T_{\mathrm{cut}}$ for
$1/T$ and \eqref{ch5:upper:eq:horizon-monotonicity} for the learning term.
After multiplying the queue bound by $\varepsilon^2=e^{-2\ell_\varepsilon}$,
the remaining expression involving $G$ is also decreasing in
$\ell_\varepsilon$. Indeed,
\begin{equation}
\begin{aligned}
 \frac{\mathrm d}{\mathrm d\ell_\varepsilon}
 \log\!\left(e^{-2\ell_\varepsilon}
            G(c+6\ell_\varepsilon)^2\right)
 &=-2+12\frac{G'(c+6\ell_\varepsilon)}{G(c+6\ell_\varepsilon)}\\
 &\overset{\text{(A)}}{\le}-2+\frac{24}{c+6\ell_\varepsilon}
 \overset{\text{(B)}}{<}0.
\end{aligned}
 \label{ch5:upper:eq:slack-monotonicity}
\end{equation}
Step (A) is \eqref{ch5:upper:eq:G-log-derivative} at
$z=c+6\ell_\varepsilon$. Step (B) uses $c=29/2>12$ and
$\ell_\varepsilon\ge0$. Multiplying
\eqref{ch5:upper:eq:long-cutoff-bound} by $\varepsilon^2$
therefore gives
\begin{equation}
\begin{aligned}
\varepsilon^2\overline Q_T
 &\le3\varepsilon+24e^{-c}\varepsilon^7
 +\frac{8192}{81}e^{-c}e^{-2\ell_\varepsilon}
   G(c+6\ell_\varepsilon)^2\\
 &\overset{\text{(A)}}{\le}3+e^{-c}
   \left(24+\frac{8192}{81}G(c)^2\right).
\end{aligned}
 \label{ch5:upper:eq:long-horizon-bound}
\end{equation}
Step (A) uses $\varepsilon\le1$ for the first two terms and
\eqref{ch5:upper:eq:slack-monotonicity} to bound the last term by its
value at $\ell_\varepsilon=0$.
For the fixed choice $c=29/2$, both numerical coefficients are smaller
than $170$:
\begin{align}
\frac78+\sqrt{\frac{128c}{9}}+\frac{32c}{3}
 &=\frac78+\frac{8\sqrt{29}}3+\frac{464}3<169.903,
\label{ch5:upper:eq:short-numerical-constant}\\
G(c)&=\frac13+8\frac{29}{2}
               +8\left(\frac{29}{2}\right)^2
             =\frac{5395}{3},\notag\\
3+e^{-c}\left(24+\frac{8192}{81}G(c)^2\right)
 &=3+e^{-29/2}
   \left(24+\frac{8192}{81}\left(\frac{5395}{3}\right)^2\right)
<167.959.
\label{ch5:upper:eq:long-numerical-constant}
\end{align}
These are direct evaluations of the constants in
\eqref{ch5:upper:eq:short-horizon-bound} and
\eqref{ch5:upper:eq:long-horizon-bound}, with the displayed decimals
rounded upward.
\endgroup
